%% ============================================================
%%  AttnFuse — Glossary of Terms
%%  Plain-English definitions of every piece of jargon used
%%  in the project, the report, and the slides.
%% ============================================================
\documentclass[11pt]{article}

\usepackage[margin=0.9in]{geometry}
\usepackage[T1]{fontenc}
\usepackage{lmodern}
\usepackage{microtype}
\usepackage{amsmath,amssymb}
\usepackage{xcolor}
\usepackage[colorlinks=true,linkcolor=blue!60!black,
            citecolor=blue!60!black,urlcolor=blue!60!black]{hyperref}
\usepackage{parskip}
\usepackage{enumitem}
\usepackage{titlesec}

\titleformat{\section}{\Large\bfseries\color{blue!60!black}}{\thesection.}{0.6em}{}
\titleformat{\subsection}{\large\bfseries\color{black!75}}{}{0pt}{}

% Definition macro: \term{Name}{plain-english definition}
\newcommand{\term}[2]{%
  \par\smallskip
  \noindent\textbf{\color{blue!50!black}#1.} #2\par
}
% Definition with a "why it matters" follow-up
\newcommand{\termw}[3]{%
  \par\smallskip
  \noindent\textbf{\color{blue!50!black}#1.} #2 \emph{\color{black!60}Why it matters here:} #3\par
}

\title{\textbf{AttnFuse Glossary}\\[4pt]
       \large Every piece of jargon in the project, in plain English}
\author{}
\date{}

\begin{document}
\maketitle

\noindent\textit{Use this as a dictionary. Each term has a one-sentence plain
explanation; for the ones that show up often in the AttnFuse code, the report,
or the slides, there is also a short ``why it matters here'' line so you can
see how the concept fits into the project.}

\tableofcontents
\bigskip

%% =============================================================
\section{AI \& Attention Basics}
%% =============================================================

\termw{Transformer}
{The neural-network architecture behind GPT, BERT, LLaMA, etc. Made of repeated
``attention'' and ``feed-forward'' blocks.}
{Every variant we accelerate lives inside a transformer.}

\termw{Self-attention}
{The step where every word in the input looks at every other word and figures
out which ones to pay attention to.}
{This is the whole thing AttnFuse optimises.}

\term{Cross-attention}
{Like self-attention, but the ``looking'' tensor (queries) comes from one
sequence and the ``looked-at'' tensor (keys/values) from another. Used in
encoder--decoder models. \emph{Not yet supported by AttnFuse.}}

\termw{Q, K, V (Query, Key, Value)}
{Three vectors derived from each word. $Q$ is ``what I'm looking for,''
$K$ is ``what I am'' (used for matching), $V$ is ``what I carry'' (the
actual information to mix in).}
{Every attention combinator in AttnFuse takes $Q$, $K$, $V$ as input.}

\termw{Attention score matrix ($S$)}
{The $N \times N$ table of numbers you get from $S = QK^\top / \sqrt{d}$.
Entry $S_{i,j}$ says ``how much should I care about word $j$ when reading
word $i$?''}
{This is the matrix we never want to write to slow GPU memory.}

\termw{Softmax}
{The function that turns raw scores into percentages that add up to 1.
$\text{softmax}(x_i) = e^{x_i} / \sum_j e^{x_j}$.}
{Turns the score matrix into attention probabilities.}

\termw{Online softmax (Milakov \& Gimelshein, 2018)}
{A way to compute softmax tile-by-tile, keeping just a running max and a
running sum, without ever holding the full row in memory.}
{This is the algorithmic trick that lets FlashAttention --- and AttnFuse ---
process attention without writing out the $N \times N$ score matrix.}

\termw{Head / multi-head attention}
{Attention is run multiple times in parallel with different learned
projections; each parallel run is a ``head.'' GPT-2 small has 12 heads.}
{The kernel processes all heads in parallel; \texttt{num\_heads} is one
of the kernel parameters.}

\termw{head\_dim ($D$)}
{The dimensionality of $Q$, $K$, $V$ per head. GPT-2 uses $D{=}64$;
LLaMA uses $D{=}128$.}
{One of the constexpr knobs in the kernel. The tile table is indexed by it.}

\term{N (sequence length)}
{The number of tokens in the input. We benchmark $N \in \{512, 1024, 2048, 4096\}$.}

\term{B (batch size)}
{How many independent sequences we process at once.}

\termw{Positional encoding}
{Extra information added so the model knows which word came first, second,
etc. Pure attention is order-blind without this.}
{ALiBi and RoPE are two flavours of positional encoding, both supported.}

%% =============================================================
\section{Attention Variants}
%% =============================================================

\termw{Dense attention}
{Every word attends to every other word. Used in BERT (encoder-only).}
{The simplest variant; matches FA2 within 15\%.}

\termw{Causal attention}
{Each word can only look at itself and the words before it. Used in
autoregressive models like GPT.}
{Implemented via a lower-triangular mask. The kernel's inner loop trims
upper-triangular blocks.}

\termw{Sliding-window attention}
{Each word only looks at the last $W$ words (or $\pm W$ for bidirectional).
Used in Mistral ($W{=}4096$) and Longformer.}
{Our biggest win vs PyTorch SDPA (38$\times$ at $N{=}4096$).}

\termw{ALiBi (Attention with Linear Biases)}
{Press et al., 2022: add a small linear penalty proportional to the
distance between words. Works without learned position embeddings.}
{One of our bias combinators (\texttt{af.alibi}).}

\termw{RoPE (Rotary Position Embedding)}
{Su et al., 2021: encode position by \emph{rotating} the $Q$ and $K$
vectors by an angle that depends on the word's position. Used in LLaMA,
Mistral, Falcon.}
{Our novel fused-kernel implementation rotates inside the inner loop ---
the project's most distinctive technical contribution.}

\termw{Additive bias attention}
{Score gets an extra bias matrix added: $S' = S + B$. The bias can be
relative-position, learned, or task-specific.}
{Supported as \texttt{af.additive\_bias} --- bias passed at call time, no
recompilation.}

\termw{Block-sparse attention (BigBird-style)}
{A sparsity pattern combining global tokens + local windows + random
attention. Lets you process very long sequences sub-quadratically.}
{\emph{Not yet implemented} in AttnFuse; mentioned in the proposal and
listed under future work.}

\termw{MQA / GQA (Multi-Query / Grouped-Query Attention)}
{Variants where the keys and values are shared across multiple query heads
to save memory. MQA: one K/V head total; GQA: a few groups.}
{Llama 3 uses GQA. Adding it is on the publication-extension list.}

%% =============================================================
\section{GPU Hardware Concepts}
%% =============================================================

\termw{GPU (Graphics Processing Unit)}
{A massively parallel processor designed for throughput, not latency.}
{The hardware we target. We benchmark on an RTX~3090.}

\termw{SM (Streaming Multiprocessor)}
{The basic ``core'' of a GPU --- a self-contained block that runs many
threads in parallel. RTX~3090 has 82 SMs.}
{Our kernel launches one block per (batch, head, query-tile); each block
runs on one SM.}

\termw{Warp}
{A group of 32 threads that execute the same instruction together. The
basic unit of GPU scheduling.}
{\texttt{num\_warps} per block is a tile-config knob; we mostly use 4 or 8.}

\termw{Thread}
{The smallest unit of work; warps are made of 32 threads.}
{We don't write per-thread code --- Triton handles that.}

\termw{HBM (High-Bandwidth Memory)}
{The big, slow off-chip memory on a GPU. RTX~3090 has 24~GB of GDDR6X at
936~GB/s.}
{Writing the $N \times N$ score matrix here is the cost FlashAttention
eliminates.}

\termw{SMEM (Shared Memory) / SRAM}
{The small, fast on-chip memory per SM. RTX~3090 has roughly 101~KB
usable per block.}
{Where Q, K, V tiles live during the tiled loop. The tile-config decision
is dominated by ``how much fits in SMEM.''}

\termw{Register file}
{Even faster on-chip storage --- thousands of 32-bit registers per warp.}
{Where the running max $m$, running sum $\ell$, and accumulator
$\mathbf{acc}$ live during online softmax.}

\termw{Tensor Cores}
{Special hardware on Ampere+ GPUs that does small matrix multiplies very
fast (in fp16/bf16). RTX~3090 peak: 142 TFLOPS fp16.}
{All our matmuls (Q$\cdot$K$^\top$, P$\cdot$V) hit tensor cores.}

\termw{TFLOPS (Teraflops per second)}
{Trillions of floating-point operations per second. The throughput
metric.}
{Our headline numbers (105 TFLOPS causal, 298 TFLOPS sliding-window).}

\termw{HBM bandwidth (GB/s)}
{How fast you can read/write off-chip memory. RTX~3090: 936~GB/s.}
{The denominator of the roofline analysis.}

\termw{Arithmetic intensity (FLOP/byte)}
{Math operations per byte of memory traffic. A kernel with low intensity
is bandwidth-bound; high intensity is compute-bound.}
{Plotted on the X-axis of the roofline figure.}

\termw{Ridge point}
{The arithmetic intensity at which a kernel transitions from
memory-bound to compute-bound. RTX~3090: $\sim$152 FLOP/byte.}
{If your kernel is below this number, optimising compute is wasted ---
fix memory traffic first.}

\termw{Roofline model (Williams et al., 2009)}
{A plot showing peak FLOPS vs peak bandwidth and where your kernel sits.}
{Used in our evaluation; \texttt{benchmarks/roofline\_runner.py}.}

\termw{Compute-bound / Memory-bound}
{A kernel is compute-bound when math is the bottleneck, memory-bound when
moving data is. Most attention kernels at long $N$ are compute-bound.}
{Helps explain why dense at small $N$ underperforms.}

\termw{Occupancy}
{Fraction of the GPU's parallel slots actually filled. Higher is usually
better, up to a point.}
{Tile-config decisions balance occupancy against per-block resource use.}

\termw{Ampere (sm\_86) / Hopper (sm\_90)}
{NVIDIA GPU generations. RTX~3090 = Ampere; A100 = Ampere data-center;
H100 = Hopper.}
{Our tile table is Ampere-tuned. Re-tuning for Hopper is on the
extension list.}

\termw{PTX}
{NVIDIA's pseudo-assembly language --- one level above actual GPU machine
code. Triton compiles down to PTX.}
{Cached by Triton in \texttt{\textasciitilde/.triton/cache/}.}

%% =============================================================
\section{Programming Models for GPUs}
%% =============================================================

\termw{Kernel}
{A function that runs on the GPU, launched from the CPU.}
{The compiled output of AttnFuse is a single kernel per variant.}

\termw{Kernel launch}
{The act of telling the GPU ``run this kernel with this grid of blocks
on these arguments.''}
{Has a fixed overhead ($\sim$5~$\mu$s). One reason we fuse multiple
operations into one launch.}

\termw{Fused kernel}
{A single kernel that does what several separate kernels would do, by
keeping intermediate values in fast memory.}
{The reason AttnFuse exists. Naive attention launches 5+ kernels; we
launch 1.}

\termw{CUDA}
{NVIDIA's low-level GPU programming language (C++ flavoured). Maximum
control, maximum verbosity. FlashAttention is written in CUDA.}
{We \emph{don't} write CUDA --- we target Triton instead.}

\termw{Triton (Tillet et al., 2019)}
{A Python-embedded language for writing GPU kernels at the tile level
rather than the thread level. Compiles to PTX via LLVM.}
{Our codegen target. Roughly 85\% of CUDA performance for $\sim$10\% of
the work.}

\termw{\texttt{@triton.jit}}
{Decorator that marks a Python function as a Triton kernel. The function
is JIT-compiled to PTX on first call.}
{All AttnFuse kernels are \texttt{@triton.jit}.}

\termw{\texttt{tl.constexpr}}
{Triton type marker that says ``this argument is known at compile time,
specialise the kernel for its value.''}
{The mechanism behind ``one kernel, many variants.'' \texttt{MASK\_KIND},
\texttt{BIAS\_KIND}, etc. are all constexpr.}

\termw{JIT (Just-In-Time) compilation}
{Compile when first called, not ahead of time. Triton uses JIT so it can
specialise per shape.}
{Cost: 0.8--1.8 seconds for our first call. Subsequent calls hit the cache.}

\termw{LLVM}
{The compiler backend Triton uses to lower its IR to PTX.}
{Most of Triton's compile time is actually LLVM optimisation.}

\termw{Autotune}
{Triton's built-in mechanism to try multiple tile configs at JIT time
and pick the fastest.}
{We \emph{don't} use this; we have a hand-picked tile table. Future work.}

%% =============================================================
\section{Compiler \& Language Concepts}
%% =============================================================

\termw{DSL (Domain-Specific Language)}
{A small language designed for a narrow problem. SQL, HTML, regex are DSLs.}
{AttnFuse is a DSL for attention.}

\termw{Embedded DSL}
{A DSL that lives inside a host language --- you use the host's syntax,
the DSL is just a library.}
{AttnFuse is embedded in Python; combinators are Python functions.}

\termw{Operator overloading}
{Redefining what \texttt{+}, \texttt{@}, etc.\ mean on your own objects.}
{We override \texttt{@} on the \texttt{TensorSym} class so
\texttt{P @ V} builds a graph node instead of multiplying.}

\termw{Symbolic tracing}
{Running a function once with placeholder objects to record what
operations it would perform, without actually doing them.}
{How we capture the user's attention function into a graph.}

\termw{IR (Intermediate Representation)}
{A data structure the compiler manipulates --- between user code and
machine code. ``Halfway'' representation.}
{AttnFuse has two: a high-level graph and a low-level TiledKernel.}

\termw{High-level IR (the Graph)}
{Captures \emph{what the user wrote} --- a DAG of typed nodes
(ScoreOp, MaskOp, BiasOp, NormOp, MatMulPV).}
{Lives in \texttt{attnfuse/ir/high\_level.py}. Hashed to make the cache key.}

\termw{Low-level IR (TiledKernel)}
{Captures \emph{how the kernel should be tiled} --- block sizes,
constexpr flags, scale factors. Flat record, no graph structure.}
{Lives in \texttt{attnfuse/ir/tiled.py}.}

\termw{Lowering}
{Translating a higher-level IR to a lower-level one, removing abstraction.}
{Our \texttt{lower\_to\_tiled} pass.}

\termw{Compiler pass}
{A single transformation that reads the IR and produces a new IR (or the
same IR, modified).}
{We have four: fuse, tile, lower, codegen.}

\termw{Fusion}
{Combining multiple operations so their intermediate results never leave
fast memory.}
{Score $\to$ mask $\to$ bias $\to$ softmax $\to$ matmul-with-V all fused
into one tiled loop.}

\termw{Tiling}
{Splitting a big loop into nested loops over small ``tiles'' that fit
in fast memory.}
{Our \texttt{choose\_tile\_config} pass picks tile dimensions.}

\termw{Codegen}
{The compiler stage that emits the target source code (Triton, in our
case) from the lowest IR.}
{Lives in \texttt{attnfuse/compiler/codegen.py}.}

\termw{Specialisation}
{Generating a separate compiled version of a function for each unique set
of compile-time inputs.}
{Triton specialises our kernel per (MASK\_KIND, BIAS\_KIND, NORM\_KIND,
ROPE\_KIND, BLOCK\_M, BLOCK\_N, HEAD\_DIM) tuple.}

\termw{Dataflow graph / DAG}
{A directed acyclic graph where nodes are operations and edges are
data dependencies.}
{Our high-level IR.}

\termw{AST (Abstract Syntax Tree)}
{The tree representation of source code after parsing. We \emph{don't}
work with Python ASTs --- we trace symbolically instead.}

\termw{MLIR}
{Multi-Level IR --- an LLVM project for building reusable compiler IRs.
Has many ``dialects'' for different domains.}
{Inspired our two-level design, but we don't depend on it. Writing an
MLIR dialect was considered and rejected as overkill for one project.}

\termw{TVM}
{An end-to-end deep-learning compiler with schedule/compute separation.}
{Comparable scope; broader than attention. We took inspiration but
implemented our own pipeline.}

\termw{Halide}
{An older image-processing DSL with schedule/compute separation; an
ancestor of TVM.}
{Conceptual prior art, not a direct comparison.}

%% =============================================================
\section{Numerical Formats}
%% =============================================================

\termw{fp32 (single precision)}
{32-bit floating point. Default for general-purpose computing.}
{Our slowest dtype on GPU; tensor cores prefer 16-bit.}

\termw{fp16 (half precision)}
{16-bit floating point. Half the memory, twice the throughput on tensor
cores. Limited range can cause underflow.}
{Our primary benchmark dtype.}

\termw{bf16 (brain float, bfloat16)}
{16-bit floating point with the same exponent range as fp32 but reduced
mantissa. Same speed as fp16, fewer underflow problems.}
{Standard for LLM training. We benchmark all three (fp16, bf16, fp32).}

\termw{fp8 / int8}
{8-bit formats used in inference quantisation. Not currently supported
by AttnFuse; future work.}

\termw{Numerical stability}
{Whether a computation gives a sensible answer despite floating-point
limits.}
{Softmax is unstable without the max-subtraction trick. Online softmax
preserves that trick across tiles.}

\termw{Max-subtraction trick}
{Compute $\text{softmax}(x)$ as $e^{x_i - m}/\sum_j e^{x_j - m}$ where
$m = \max(x)$. Same answer, no overflow.}
{The reason online softmax works.}

%% =============================================================
\section{AttnFuse-Specific Vocabulary}
%% =============================================================

\termw{Combinator}
{A composable building block --- in AttnFuse, a function that takes
graph nodes and returns a graph node.}
{We have seven: \texttt{scaled\_dot\_product}, \texttt{rope},
\texttt{causal}, \texttt{sliding\_window}, \texttt{alibi},
\texttt{additive\_bias}, \texttt{softmax}.}

\termw{\texttt{@af.attention}}
{The decorator the user puts on their attention function. It triggers
tracing, lowering, codegen, and caching on first call.}

\termw{\texttt{TensorSym}}
{Symbolic placeholder for a tensor. Doesn't carry data --- only shape,
dtype, and a name. Used during tracing.}

\termw{ScoreOp / MaskOp / BiasOp / NormOp / MatMulPV}
{Node types in the high-level IR. Each combinator emits one of these.}

\termw{MASK\_KIND}
{Integer constexpr in the kernel: $0 =$ none, $1 =$ causal, $2 =$ sliding-
window. Selected by the user via the mask combinator.}

\termw{BIAS\_KIND}
{Integer constexpr: $0 =$ none, $1 =$ ALiBi, $2 =$ additive (external tensor).}

\termw{NORM\_KIND}
{Integer constexpr: $0 =$ softmax, $1 =$ ReLU attention.}

\termw{ROPE\_KIND}
{Integer constexpr: $0 =$ off, $1 =$ fused RoPE on.}

\termw{BLOCK\_M, BLOCK\_N}
{The query-tile and key-tile sizes in the inner loop. Typically 128 and
64 for fp16 on Ampere.}

\termw{num\_warps}
{Warps per block (Triton parameter). Usually 4 or 8.}

\termw{num\_stages}
{Pipeline depth for asynchronous memory loads. Higher = more overlap of
load and compute, more SMEM use.}
{We use 3 for fp16, 2 for fp32 and for RoPE variants (SMEM-constrained).}

\termw{LaunchBundle}
{Cache entry holding everything needed to launch a compiled kernel:
the JIT-compiled function, the constexpr dictionary, placeholder
slope/bias/RoPE tensors, and the launch grid.}
{Reduces per-call dispatch overhead from $\sim$7~$\mu$s to $<$1~$\mu$s.}

\termw{Dispatch}
{The runtime step that takes user arguments, looks up the cached kernel,
and launches it.}
{Lives in \texttt{attnfuse/runtime/dispatch.py}.}

\termw{Graph signature}
{A SHA-1 hash of the high-level IR's structural content. The cache key.}
{Same graph $\Rightarrow$ same signature $\Rightarrow$ same compiled kernel.}

\termw{\_generated/}
{Directory where we write generated Triton kernel source. Lets Triton's
JIT read the source via \texttt{importlib} (needed for LLVM hashing).}

\termw{Placeholder tensors}
{Tiny $(1,1)$ tensors used for unused kernel arguments (e.g., the bias
pointer when BIAS\_KIND${=}0$). Cached via \texttt{lru\_cache}.}

%% =============================================================
\section{Benchmarking \& Evaluation Terms}
%% =============================================================

\termw{Latency}
{Time from launch to result, usually in milliseconds.}
{Our primary speed metric.}

\termw{Throughput}
{Work done per unit time, usually TFLOPS for kernels or tokens/second
for models.}

\termw{Wall-clock time}
{Real elapsed time, as opposed to CPU time. Synonymous with latency for
GPU work.}

\termw{CUDA events}
{Hardware timestamps that measure GPU work accurately, unaffected by
CPU--GPU sync overhead.}
{We use \texttt{torch.cuda.Event} for all timings.}

\termw{Warmup}
{Initial kernel runs that we discard from timing, to let caches warm up
and JIT compile.}
{5 warmup + 50 timed iterations, median reported.}

\termw{Baseline}
{The thing you compare against. We use three: \texttt{naive} (eager
PyTorch), \texttt{sdpa} (PyTorch's flash backend), and \texttt{flash\_ref}
(our hand-written Triton FA2).}

\termw{\texttt{F.scaled\_dot\_product\_attention} (SDPA)}
{PyTorch's official attention call. Dispatches to a fixed set of
backends including FlashAttention-2.}
{Our main comparison target. SDPA falls back to slow code for
non-standard masks.}

\termw{FlashAttention / FA2 (Dao 2022, 2023)}
{The hand-written CUDA kernel that pioneered IO-aware tiled attention.
What SDPA dispatches to under the hood.}
{Our performance ceiling for the dense and causal variants.}

\termw{Speedup}
{Ratio of baseline time to our time. $38\times$ means we are 38 times
faster.}

\termw{Ablation study}
{Sweeping one parameter at a time to see its effect.}
{Our tile-config study sweeps BLOCK\_M, BLOCK\_N, num\_warps, num\_stages.}

\termw{Peak memory}
{The highest GPU memory usage during the run, in MB.}
{The ``thrifty'' metric --- what makes long contexts feasible.}

\termw{Nsight Compute (\texttt{ncu})}
{NVIDIA's kernel-level profiler. Measures real tensor-core utilisation,
real HBM traffic, occupancy.}
{Not available on our evaluation host --- our TC\% and HBM~GB/s are
computed analytically.}

%% =============================================================
\section{Tools \& Ecosystem}
%% =============================================================

\termw{PyTorch}
{The deep-learning framework we use as a baseline and dependency.}

\termw{\texttt{torch.compile}}
{PyTorch's recent graph-compiler backend (TorchDynamo + Inductor).
Auto-fuses simple op patterns.}
{Adjacent technology; we don't go through it.}

\termw{\texttt{flex\_attention} (PyTorch 2.5+, 2024)}
{A newer PyTorch API where you supply a \texttt{score\_mod} function to
modify attention scores; it generates a fused kernel automatically.}
{Our closest competitor. Cannot fuse RoPE (which acts on Q, K \emph{before}
the matmul, not on the score).}

\termw{xFormers}
{Meta's library of hand-written attention kernels.}
{Adjacent prior art; not a DSL.}

\termw{Liger Kernel}
{A library of fused Triton implementations of common LLM operators.}
{Adjacent; not a DSL.}

\termw{HuggingFace transformers}
{The standard library of pretrained transformer models.}
{Not currently integrated; a future paper extension.}

\termw{pytest}
{The Python testing framework. Our 25 correctness tests run under it.}

%% =============================================================
\section{Quick-Reference: Acronyms at a Glance}
%% =============================================================

\begin{itemize}[leftmargin=1.4em,itemsep=1pt]
  \item \textbf{ALiBi} --- Attention with Linear Biases (positional bias).
  \item \textbf{AST} --- Abstract Syntax Tree.
  \item \textbf{BERT} --- Bidirectional Encoder Representations from Transformers.
  \item \textbf{bf16} --- Brain float, 16-bit.
  \item \textbf{CUDA} --- NVIDIA's GPU programming language.
  \item \textbf{DAG} --- Directed Acyclic Graph.
  \item \textbf{DSL} --- Domain-Specific Language.
  \item \textbf{FA2} --- FlashAttention-2.
  \item \textbf{fp16 / fp32 / fp8} --- Floating-point, 16 / 32 / 8 bit.
  \item \textbf{GQA / MQA} --- Grouped-Query / Multi-Query Attention.
  \item \textbf{HBM} --- High-Bandwidth Memory (GPU's main memory).
  \item \textbf{IR} --- Intermediate Representation.
  \item \textbf{JIT} --- Just-In-Time (compilation).
  \item \textbf{LLM} --- Large Language Model.
  \item \textbf{MLIR} --- Multi-Level Intermediate Representation.
  \item \textbf{MLSys} --- Machine Learning and Systems conference.
  \item \textbf{NCU} --- Nsight Compute (NVIDIA's profiler).
  \item \textbf{PTX} --- NVIDIA's pseudo-assembly (below CUDA).
  \item \textbf{QKV} --- Query, Key, Value (the three attention projections).
  \item \textbf{RoPE} --- Rotary Position Embedding.
  \item \textbf{SDPA} --- Scaled Dot-Product Attention (PyTorch's API).
  \item \textbf{SMEM} --- Shared Memory (GPU's fast on-chip memory).
  \item \textbf{SM} --- Streaming Multiprocessor.
  \item \textbf{sm\_86 / sm\_90} --- Compute capability (Ampere / Hopper).
  \item \textbf{SW} --- Sliding-Window (attention).
  \item \textbf{TFLOPS} --- Trillions of FLoating-point OPerations per Second.
  \item \textbf{TVM} --- Apache TVM, a deep-learning compiler.
\end{itemize}

\end{document}
